% =====================================================================
% Rapport final — sujet 39, The Blind Spot Paradox
%
% Gabarit IEEEtran conference, deux colonnes : celui du manuscrit critiqué
% (articleA_blindspot_v64_camera_ready.tex). Contrainte de l'énoncé :
% 10 à 15 pages, 15 étant un maximum strict.
%
% État : question A intégrée (condensée depuis redaction_QA_course_concurrente).
%        B, C, D à reprendre — voir les marqueurs en fin de fichier.
%
% Deux colonnes : le matériel large (encadré de notations, tables à plus de
% quatre colonnes) passe en table*, qui s'étale sur les deux colonnes.
%
% Compiler depuis 04_experimentations/.
% =====================================================================

\documentclass[10pt,conference]{IEEEtran}

\usepackage[utf8]{inputenc}
\usepackage[T1]{fontenc}
\usepackage[english]{babel}
\usepackage{amsmath,amssymb,amsthm}
\usepackage{graphicx}
\usepackage{booktabs}
\usepackage{microtype}
\usepackage[hidelinks,colorlinks=true,linkcolor=black,citecolor=blue,urlcolor=blue]{hyperref}

\newtheorem{theorem}{Theorem}
\newtheorem{proposition}[theorem]{Proposition}
\newtheorem{corollary}[theorem]{Corollary}
\theoremstyle{definition}
\newtheorem{remark}[theorem]{Remark}

\IEEEoverridecommandlockouts
\title{The Blind Spot Paradox: Formalizing the Race\\and Measuring the Adaptation}
\author{\IEEEauthorblockN{Alexandre Boyer, Melaine Gouillou,
Salomé Fonvielle, Ulysse Petit-Tichanné}
\IEEEauthorblockA{Filière Recherche --- Subject 39}}

\begin{document}
\onecolumn

% ---------------------------------------------------------------------
% Page de garde
% ---------------------------------------------------------------------
\thispagestyle{empty}
\vspace*{\fill}
\begin{center}
  {\LARGE\bfseries The Blind Spot Paradox\\[0.4em]
   Formalizing the Race and Measuring the Adaptation\par}
  \vspace{2.5em}
  {\large Alexandre Boyer \quad Melaine Gouillou\\[0.4em]
   Salomé Fonvielle \quad Ulysse Petit-Tichanné\par}
  \vspace{2em}
  {Filière Recherche --- Subject 39\par}
  \vspace{1em}
  % TODO : date de remise
\end{center}
\vspace*{\fill}
\newpage

% ---------------------------------------------------------------------
% Mise en contexte et notations — pleine page, avant le corps
% ---------------------------------------------------------------------
\section*{Context and Scope}

% TODO : mise en contexte (une page avec les notations). Doit poser le
% phénomène, les deux améliorations demandées par les relecteurs, et
% annoncer le partage théorique (A, B) / expérimental (C, D).
Placeholder.

\section*{Notation Used Throughout}

\begin{center}
\begin{tabular}{@{}ll@{}}
  \toprule
  $\tau^*$ & the drift instant, known by construction \\
  $\tau_{\mathrm{ARF}}$ & delay until the first tree is replaced, $\min_i \tau_i$ over the $M$ trees \\
  $\tau_{\mathrm{det}}$ & delay until the first external alarm, possibly $+\infty$ \\
  $F_A$, $F_D$ & cumulative distribution functions of $\tau_{\mathrm{ARF}}$ and $\tau_{\mathrm{det}}$ \\
  $H$ & post-drift observation horizon, $H = 2\,000$ \\
  $\lambda$, $\delta_P$ & the external detector's threshold and tolerance \\
  $\Delta e$ & drift amplitude, the error jump a non-adapting model would suffer \\
  $\mathrm{Miss}$ & the ``blind spot'' event, defined in Section~\ref{sec:miss} \\
  \midrule
  \multicolumn{2}{@{}l@{}}{\emph{for the inter-tree analysis:}} \\
  $M = 10$ & number of trees in the forest \\
  $\tau_i$ & adaptation delay of tree $i$, so $\tau_{\mathrm{ARF}} = \min_i \tau_i$ \\
  $S$ & the complete realization of the stream (one seed) \\
  $U_i$ & the randomness specific to tree $i$, identified in Section~\ref{sec:factorization} \\
  $F$ & CDF of a \emph{single} tree $\tau_i$ --- distinct from $F_A$, which is that of $\tau_{\mathrm{ARF}}$ \\
  $G_S(s)$ & $:= P(\tau_1 > s \mid S)$, the single-tree survival function \emph{given} the stream \\
  $\beta$ & reliability target, the share of drifts the monitor should catch ($0.95$ in the manuscript) \\
  \bottomrule
\end{tabular}
\end{center}

\noindent Both delays are counted from $\tau^*$ and take integer values, a fact the
Fréchet bounds of Question~A rely on. Every measured figure in this report comes from the
instrumented campaign of $20$ amplitudes $\times\ 100$ seeds $= 2\,000$ runs,
$H = 2\,000$, documented in \texttt{resultats\_R2/}.

\newpage
\twocolumn

% =====================================================================
\section{The Miss Probability: Definition, Identification, Bounds}
% =====================================================================

We formalize and bound $P_{\mathrm{miss}} = P(\tau_{\mathrm{ARF}} <
\tau_{\mathrm{det}})$ without any assumption on the dependence between the two clocks.

\subsection{Defining the Miss Event}
\label{sec:miss}

Everything that follows is a statement about $P_{\mathrm{miss}}$, and none of it means
anything until $\mathrm{Miss}$ is unambiguous. The natural reading --- the forest repairs
itself before the alarm sounds --- leaves three cases open: the two clocks coincide, the
detector never fires, or neither mechanism ever acts. We settle them by one definition
rather than three rules.

\paragraph{1. A strict inequality.} Both clocks count integer steps, so they can land on
the same one and $P(\tau_{\mathrm{ARF}} = \tau_{\mathrm{det}})$ can be strictly positive:
the choice is not cosmetic, and no continuity argument disposes of it --- real-valued
times can coincide too, all the more so for two functionals of the same error stream.
It must therefore be settled on physical grounds. A tie means the alarm sounds
\emph{exactly} when the forest repairs itself: the operator is informed at that instant
and nothing is lost, unlike the blind spot proper where the alarm comes late or never. A
tie is a detection success, whence
\[
  \mathrm{Miss} \;:=\; \{\tau_{\mathrm{ARF}} < \tau_{\mathrm{det}}\}
\]
on the extended reals, with strict inequality.
The convention carries almost no weight in practice: ties occur $0$ times out of
$1\,906$ observed races at $\lambda = 8$, $6$ out of $778$ at $\lambda = 25$, $0$ out of
$1$ at $\lambda = 50$, so the non-strict reading would move $P_{\mathrm{miss}}$ by at
most $0.003$. They are rare not for any reason of continuity but because the two
mechanisms operate on visibly different time scales --- which is the substance of the
paradox rather than a technicality.

\paragraph{2. A silent repair is a miss.} On
$\{\tau_{\mathrm{ARF}} < +\infty,\ \tau_{\mathrm{det}} = +\infty\}$ the forest repaired
itself and the alarm never sounded --- the blind spot in its purest form, worse than
$\tau_{\mathrm{ARF}} < \tau_{\mathrm{det}} < +\infty$ where the operator is at least
informed late. A finite $\tau_{\mathrm{ARF}}$ is smaller than $+\infty$, so the event
falls inside $\mathrm{Miss}$ on its own.

It must. Excluding it would restrict the count to runs where the detector eventually
spoke --- the very selection bias this question holds against the manuscript, since those
discarded are the runs where it failed most completely. At $\lambda = 50$, where $1\,999$
runs out of $2\,000$ never fire, the rate would be computed over a single run.

\paragraph{3. A double starvation is not.} On
$\{\tau_{\mathrm{ARF}} = +\infty,\ \tau_{\mathrm{det}} = +\infty\}$ neither mechanism
ever acted. The paradox rests on a \emph{silent repair} erasing the evidence; if the
repair never happened there is nothing to erase, and the model simply stays broken. Since
$\infty < \infty$ is false, the event is excluded on its own. Its empirical frequency is
$0$ at all three thresholds, so the convention is never exercised on this grid --- but
nothing in the model forbids it, and a slower forest would produce it.

\begin{center}
\fbox{\parbox{0.92\columnwidth}{\textbf{Adopted convention.}
\[ \mathrm{Miss} := \{\tau_{\mathrm{ARF}} < \tau_{\mathrm{det}}\} \qquad \text{(strict, extended reals)} \]
A tie goes to the detector, a silent repair is a miss, a double starvation is not --- all
three from the definition, none from a rule of its own. That is what makes it
trustworthy rather than fitted to the answer.}}
\end{center}

\subsection{Making the Miss Probability Observable}

$\mathrm{Miss}$ is defined on the complete race, with no time limit; the campaign
observes $H = 2\,000$ steps. On runs where neither clock has fired by then the winner is
not known, and $P_{\mathrm{miss}}$ depends on what was never seen: it is not, as it
stands, computable from the data. Two ways out are unacceptable --- assuming that
censored runs behave like the others cannot be checked, and discarding them is the
selection bias just rejected. The aim is therefore to replace $P_{\mathrm{miss}}$ by the
probability of an event that \emph{is} observable, or failing that to trap it between two
such probabilities.

\paragraph{What the horizon lets us see.} Absolute censoring is
$C_H := \{\min(\tau_{\mathrm{ARF}},\tau_{\mathrm{det}}) > H\}$, on which the race is
undecided. Sorting runs by the two observable events $\{\tau_{\mathrm{ARF}} \le H\}$ and
$\{\tau_{\mathrm{det}} \le H\}$ gives a partition where the verdict is known everywhere
but one cell: regions~1 ($\le H$, $\le H$, decided directly) and~2 ($\le H$, $>H$, a
certain Miss) together form exactly $\{\tau_{\mathrm{ARF}} \le H\} \cap \mathrm{Miss}$;
region~3 ($>H$, $\le H$) is never a Miss; region~4 is $C_H$, undetermined.

\paragraph{Trapping the unknown between two computable numbers.} Region~4 can shift the
total by its own mass and no more. Write
$L := P(\{\tau_{\mathrm{ARF}} \le H\} \cap \mathrm{Miss})$, computable from the data with
no assumption about what happens beyond $H$. Counting region~4 as never a Miss gives the
lower bound, entirely as Miss the upper one:
\[
  L \;\le\; P_{\mathrm{miss}} \;\le\; L + P(C_H).
\]
This is a Manski-type bound: nothing is assumed about the unobserved
runs, their possible influence is bounded by how many there are, and the width of the interval is exactly
$P(C_H)$ --- the extent of our ignorance.

\paragraph{On this campaign, the interval collapses.} $\tau_{\mathrm{ARF}}$ is
\emph{never} censored: $0$ runs out of $2\,000$, the largest value observed being
$1\,293$ steps. The forest, at maximal reactivity ($c_{\mathrm{int}} = 1$), always
replaces a tree before the horizon, so $\min(\tau_{\mathrm{ARF}},\tau_{\mathrm{det}})$
is always finite, $C_H = \varnothing$, and the two bounds coincide:
\[
  P_{\mathrm{miss}} \;=\; P(\{\tau_{\mathrm{ARF}} \le H\} \cap \mathrm{Miss}).
\]

\begin{center}
\fbox{\parbox{0.92\columnwidth}{\textbf{Result.} The unconditional miss probability equals
the probability of an event observable within the horizon. This is an equality, not an
assumption: the naive computation would have returned the same number, but nothing would
have told us it was right. What the construction adds is the guarantee --- and the
condition under which it would fail. Were the forest able to stay inert past $H$,
$P_{\mathrm{miss}}$ would cease to be identified and only the interval would survive.}}
\end{center}

\subsection{Universal Fréchet Bound}

We now drop the finite horizon and ask what the two marginal laws concede on their own,
with no assumption on the coupling. This is a different kind of ignorance from the one
just treated: there, data were missing; here, a structure is.

\paragraph{Lower bound.} For a fixed $s \in \mathbb{N}$, the event
$\{\tau_{\mathrm{ARF}}\le s < \tau_{\mathrm{det}}\}$ implies
$\tau_{\mathrm{ARF}}<\tau_{\mathrm{det}}$, so it is a \emph{subset} of $\mathrm{Miss}$.
Write $A = \{\tau_{\mathrm{ARF}} \le s\}$ and $B = \{\tau_{\mathrm{det}} > s\}$, so
$P(A) = F_A(s)$ and $P(B) = 1 - F_D(s)$. Fréchet's inequality follows from
$P(A \cup B) \le 1$:
\[
  P(A \cap B) \;=\; P(A) + P(B) - P(A \cup B) \;\ge\; P(A) + P(B) - 1 ,
\]
and since a probability is non-negative,
$P(A \cap B) \ge \max\big(0, P(A)+P(B)-1\big)$. Substituting,
$P(A) + P(B) - 1 = F_A(s) - F_D(s)$. The inclusion holding for every $s$, we keep the
largest value:
\[
  P_{\mathrm{miss}} \;\ge\; \max\Big(0,\ \sup_{s\in\mathbb{N}} \big[F_A(s)-F_D(s)\big]\Big).
\]

\paragraph{Upper bound.} Here we proceed by \emph{superset}. We claim
\[
  \mathrm{Miss} \;\subseteq\; \{\tau_{\mathrm{ARF}}\le s\} \cup \{\tau_{\mathrm{det}}> s+1\},
\]
and prove it by taking a run in $\mathrm{Miss}$ and splitting on where
$\tau_{\mathrm{ARF}}$ falls. If $\tau_{\mathrm{ARF}} \le s$ the run lies in the first set
and nothing more is needed. If $\tau_{\mathrm{ARF}} \ge s+1$, then
$\tau_{\mathrm{det}} > \tau_{\mathrm{ARF}} \ge s+1$ puts it in the second. In words:
either the forest was early, or the alarm --- which comes after it --- was late. Boole's
subadditivity, then optimizing over $s$ and truncating at $1$, gives
\[
  P_{\mathrm{miss}} \;\le\; \min\Big(1,\ \inf_{s\in\mathbb{N}} \big[F_A(s) + 1 - F_D(s+1)\big]\Big).
\]

\begin{center}
\fbox{\parbox{0.92\columnwidth}{\textbf{Result.} The two bounds above hold for
\emph{any} dependence structure between $\tau_{\mathrm{ARF}}$ and
$\tau_{\mathrm{det}}$: they use the marginal laws alone, and assume nothing about how
the two clocks are coupled. That is what makes them \emph{universal}.}}
\end{center}

\paragraph{Verification on a toy case.} Take both clocks uniform on $\{1,2\}$, so
$F_A = F_D$ with $F_A(1) = \tfrac12$. Then $\mathrm{Miss}$ requires
$\tau_{\mathrm{ARF}} = 1$ and $\tau_{\mathrm{det}} = 2$, and its probability depends
entirely on the coupling: $0$ when the two clocks always agree, $0.25$ under
independence, $0.5$ when one is $1$ exactly as the other is $2$. Identical marginals,
three answers. The formulas give a lower bound of $0$ and an upper bound of $0.5$
(attained at $s = 0$ and $s = 1$), both reached: the interval is exactly the set of
achievable values, and nothing sharper can be had from marginals alone. Assuming
independence, as the manuscript does, is therefore not a harmless convenience --- it
selects one point out of $[0,\,0.5]$.

\paragraph{On the campaign's marginals.} Within a cell, $\hat F_A(s)$ is the share of the
$100$ seeds with $\tau_{\mathrm{ARF}} \le s$ and $\hat F_D(s)$ the share with
$\tau_{\mathrm{det}} \le s$; the two bounds are the formulas above evaluated on these.
$\hat F_D$ is sub-stochastic --- a censored $\tau_{\mathrm{det}}$ enters it at no finite
$s$ --- so censored runs count as detector losses rather than as missing data. At
$\lambda = 8$:

\begin{table}[t]
\centering
\begin{tabular}{crrr}
\toprule
$\Delta e$ & lower & $\hat P_{\mathrm{miss}}$ & upper \\
\midrule
$0.028$ & $0.91$ & $1.00$ & $1.00$ \\
$0.141$ & $0.10$ & $0.12$ & $0.26$ \\
$0.498$ & $0.01$ & $0.01$ & $0.18$ \\
\bottomrule
\end{tabular}
\caption{Fréchet bounds on the estimated marginals, $\lambda = 8$.}
\end{table}

\noindent The bound holds on all $60$ cells of the grid
(\texttt{resultats\_R2/scripts/chiffres\_QA.py}), and is not vacuous: mean width $0.195$
at $\lambda = 8$ and $0.106$ at $\lambda = 25$. At $\lambda = 50$ it falls to $0.001$,
but only because the detector never alarms there, $\hat F_D$ stays flat at $0$, and both
formulas reach $P_{\mathrm{miss}} = 1$ on their own --- exactness by degeneracy, not by
sharpness. At the weakest amplitude the estimated lower bound of $0.91$ becomes $0.88$ at
the $95\%$ level after a bootstrap resampling the $100$ seeds. There, the blind spot is
forced by the marginals alone, with no assumption on the coupling: it is not an artefact
of an independence hypothesis.

\subsection{The Race, Measured}

The event is now defined and its probability identified, so the campaign can be read
directly. Table~\ref{tab:race} gives the partition and $P_{\mathrm{miss}}$ at the three
thresholds.

\begin{table}[t]
\centering
\begin{tabular}{crrrrl}
\toprule
$\lambda$ & region 1 & region 2 & $C_H$ & $\hat P_{\mathrm{miss}}$ & Wilson $95\%$ \\
\midrule
$8$  & $1{,}906$ & $94$    & $0$ & $0.0905$ & $[0.0787;0.1039]$ \\
$25$ & $778$     & $1{,}222$ & $0$ & $0.9715$ & $[0.9633;0.9779]$ \\
$50$ & $1$       & $1{,}999$ & $0$ & $1.0000$ & $[0.9981;1.0000]$ \\
\bottomrule
\end{tabular}
\caption{The race, measured on the campaign. Regions 3 and 4 are empty, so
$\hat P_{\mathrm{miss}}$ estimates $P_{\mathrm{miss}}$ itself rather than a bound on it.}
\label{tab:race}
\end{table}

These totals pool cells of very unequal difficulty, and the per-cell breakdown belongs to
Question~D.

% =====================================================================
\section{The Independence Assumption Between Trees}
% =====================================================================

The $M$ trees share the same stream, yet the manuscript treats the $\tau_i$ as
independent when computing $\tau_{\mathrm{ARF}} = \min_i \tau_i$. In which direction does
that bias its formula?

\subsection{Conditional Factorization}
\label{sec:factorization}

Once $S$ is fixed --- the entire trajectory $(x_t,y_t)$ known --- what still distinguishes
one tree from another is $U_i$: the Poisson weights drawn online for bagging, and the
random feature subspaces tested at each node. The first is the mechanism the manuscript
itself invokes, ``Poisson-weighted bootstrap (Oza bagging with $\lambda = 6$)''
(Section~III-G, citing Gomes et al.\ 2017); the second is a defining property of the
Adaptive Random Forest in the same reference\footnote{On the R2 stream it contributes
little: with two features only, a random subspace selects roughly one of two at each node.
Diversity there rests almost entirely on the Poisson weights, which suffices --- only the
i.i.d.\ character of the $U_i$ is needed.}.

These draws are independent across trees by construction: tree $i$ never sees tree $j$'s
weights, and a tree created at replacement inherits neither weights nor structure from the
one it replaces. Since $\tau_i$ is a deterministic function of $(S, U_i)$, fixing $S$
leaves only the i.i.d.\ variation of $U_i$. Conditionally on $S$, the $\tau_i$ are
therefore independent and identically distributed, and
\[
  P(\tau_{\mathrm{ARF}}>s \mid S) = \prod_{i=1}^{M} P(\tau_i>s\mid S) = G_S(s)^M .
\]
The independence assumption is thus not wrong in kind --- it is stated at the wrong level.
It holds given the stream; it fails once the stream is averaged out.

\subsection{The Direction of the Bias}
\label{sec:convexity}

Unconditioning by the law of total probability gives
$P(\tau_{\mathrm{ARF}}>s) = \mathbb{E}_S[G_S(s)^M]$. Since $x \mapsto x^M$ is convex on
$[0,1]$, Jensen's inequality yields
\[
  \mathbb{E}_S\big[G_S(s)^M\big] \;\ge\; \big(\mathbb{E}_S[G_S(s)]\big)^M
  = \big(1-F(s)\big)^M ,
\]
the last equality because $\mathbb{E}_S[G_S(s)] = P(\tau_1>s) = 1-F(s)$.

\begin{center}
\fbox{\parbox{0.92\columnwidth}{\textbf{Result.}
\[
  P(\tau_{\mathrm{ARF}}\le s) \;\le\; 1-\big(1-F(s)\big)^M
\]
The probability of having already adapted by time $s$ is at most what the independence
formula predicts: the manuscript \textbf{overestimates the repair speed}.}}
\end{center}

\paragraph{Equality condition.} Jensen is an equality exactly when $G_S(s)$ is almost
surely constant --- when the difficulty of adapting does not depend on which stream was
drawn. That never happens here: each seed produces a different noise trajectory after the
drift. Independence is therefore never exact, only an approximation, and a systematically
optimistic one.

\subsection{Operational Consequence}

Two quantities are affected in opposite directions, and the words below refer to different
subjects --- which is what makes the pair look contradictory. The independence formula
predicts a \emph{faster} repair than the real one, and therefore a \emph{larger}
blind-spot risk: a slower repair leaves the monitor more time, so the forest wins the race
less often than claimed. Read alone, our correction makes the blind spot \emph{less} severe
than the article reports.

\paragraph{Consequence for $M_{\mathrm{crit}}$.} Corollary~10 of the manuscript asks how
many trees the forest may carry before the monitor stops being reliable. Requiring
$P_{\mathrm{miss}} = 1-[1-F(\tau^*_{\mathrm{det}})]^M \le 1-\beta$ gives $[1-F]^M \ge
\beta$, hence $M\ln(1-F) \ge \ln\beta$; since $\ln(1-F)<0$ the division flips the
inequality:
\[
  M \;\le\; M_{\mathrm{crit}} := \left\lfloor \frac{\ln\beta}{\ln[1-F(\tau^*_{\mathrm{det}})]} \right\rfloor .
\]
$M_{\mathrm{crit}}$ is a \emph{ceiling} on the number of trees, not a floor --- the
physically sensible reading, since more trees mean a faster repair and a worse position
for the detector. As the true CDF lies below the independence one, the true
$P_{\mathrm{miss}}$ is smaller at every $M$: whatever ensemble size the certificate
admits, the real forest satisfies the guarantee too. The corollary is not invalidated but
\textbf{reclassified as a conservative ceiling}.

\paragraph{Two reservations.} The comparison holds at a \emph{fixed}
$\tau^*_{\mathrm{det}}$, which Proposition~9 makes a constant. Read as a claim about the
real blind-spot risk it assumes more, since the real $\tau_{\mathrm{det}}$ is computed on
the ensemble's error stream and so depends on how correlated the trees are. Both effects
in fact lower $P_{\mathrm{miss}}$ --- a slower repair alarms earlier, and weaker variance
reduction leaves a burstier stream that breaches $\lambda$ on noise alone --- so the sign
holds, but part of the decrease is bought with alarms that win the race without detecting
anything.

And the corollary is not rescued. At $\Delta e = 0.33$ the manuscript measures
$\widehat F(\tau^*_{\mathrm{det}}) = 0.49$: one tree alone already gives
$P_{\mathrm{miss}} = 0.49$, ten times the $5\%$ that $\beta = 0.95$ allows, whence
$M_{\mathrm{crit}} = 0$. Our correction states the true ceiling is $\ge 0$, which says
nothing --- rescuing the certificate would require $\ge 1$. The status changes, the verdict
does not.

\subsection{The Correlation Trap}

Could the dependence run the other way and reverse the conclusion? On the toy case of
two trees uniform on $\{1,2\}$, an anti-monotone coupling ($\tau_1 = 1 \iff \tau_2 = 2$)
gives $P(\tau_{\mathrm{ARF}} \le 1) = 1$ against $0.75$ under independence --- a
\emph{faster} adaptation. Its covariance is
$0 - \tfrac12\cdot\tfrac12 = -\tfrac14$, strictly negative.

No ARF can realize it. By the law of total covariance, with
$X=\mathbf{1}\{\tau_1{>}1\}$ and $Y=\mathbf{1}\{\tau_2{>}1\}$, the conditional term
vanishes by Section~\ref{sec:factorization} and both conditional expectations equal
$G_S(1)$, so
\[
  \mathrm{Cov}(X,Y) = \mathrm{Var}\big(G_S(1)\big) \;\ge\; 0 .
\]

\begin{center}
\fbox{\parbox{0.92\columnwidth}{\textbf{The trap defused.} The only channel of dependence
between two trees is the stream they share, and sharing a stream can only make them fail
together --- never in opposition. The covariance is therefore never negative, and the
anti-monotone coupling is structurally out of reach. This is not a missing link in
Section~\ref{sec:convexity}, whose proof never invokes the sign of the dependence: it is a
second consequence of the same conditional factorization, and it closes the one objection
a reader would raise against it.}}
\end{center}

% =====================================================================
% Convention de titre : les sections portent leur contenu, pas le numéro de
% la question (« Question A » a été retiré). Le correcteur doit donc trouver
% la correspondance question -> section dans « Context and Scope ».
%
%
% Phrase d'ouverture prête à l'emploi pour B, déplacée depuis la fin de A
% (elle y occupait une sous-section pour trois lignes de pont) :
%
%   With $P_{\mathrm{miss}}$ point identified, the Fréchet interval of
%   Question~A measures what the \emph{marginals alone} concede --- between
%   $0.09$ and $0.34$ of probability mass at $\lambda = 8$, depending on the
%   amplitude. Closing that gap is not a matter of more data: it requires the
%   \emph{structure} of the dependence, which is the object of this question.
% TODO : Question C --- reprendre QC1, QC2, QC3 (trois fichiers à fondre)
% TODO : Question D --- reprendre QD1 à QD4 (quatre fichiers à fondre)
% TODO : conclusion commune, et la déclaration d'usage des LLM
%
% Budget : 15 pages maximum, tout compris. La question A occupe environ
% 4 pages sous cette forme ; à surveiller quand B, C et D arriveront.
% =====================================================================

\end{document}
